\chapter{The Fences}
% Covers: outline F1-F5. Shape per turn-528 ruling: THREE sections --
% S1 "The impossibility" (F1); S2 "The lines" (F2 electron-only + F3 one honest
% paragraph referencing Ch5's standing fence, never re-entering); S3 "The fates and
% the bracket" (F4+F5, METHOD rendered as method; GAUSS lands as the MOVE fate's
% worked example -- the demotion story). Chapter language (per Ch1's arrow plan):
% the two-derivative distinction.

\section{The impossibility}
% Node-map: F1 whole -- the proved fence: a BOX-READING (inspecting finished states
% without paying elaboration) cannot separate the first variation from the second,
% therefore cannot derive the coupling; THE CLAIM STAYS NARROW (this method cannot --
% never "no method can"); choice-free stated; Frechet & Gateaux's row lands on the
% record's actual distinction (the two derivatives indistinguishable to a box-reader);
% the VINDICATION (the walk was necessary -- measurement forced, not chosen); the
% fence DISPLAYED. Gate: turn 528 (Ch7 S1, ~950 words).

Chapter 1 asked the reader to treat every fence in this book as load-bearing, and to
be suspicious of any summary that omits them. This chapter is where the fences live,
and it opens with the argument's proudest possession --- a fence that is not a
policy, not a scruple, not a boundary chosen in good taste, but a theorem.

First, the temptation it kills, stated as seductively as its victims feel it. When
the machine finishes a run, everything it did leaves remains: finished forms,
intermediate states, the whole contents of the box the run filled. All of it simply
sits there, free to inspect. A box-reading is exactly that inspection: read off the
states, compare the outputs, tabulate what stands in the box --- without ever paying
elaboration, without running anything, information after the fact and free of
charge. And the temptation follows at once: everything the machine did is recorded
in what it left behind; surely the coupling is in there somewhere, waiting to be
read rather than walked for. Why count heartbeats, measure slips, run mediant walks
to a fencepost --- why do any of this book's labor --- if a sufficiently clever
reading of the box would hand over the number?

To say why not, the book must make one distinction finer than Chapter 1 needed, and
it is the distinction on which the whole fence turns. The first difference of the
count, Chapter 1 said, is a rate; the second is a curvature; the coupling lives in
the curvature. But "curvature" hides a subtlety that two French mathematicians made
precise. A derivative can be probed one direction at a time --- displace this way,
read the response; displace that way, read again --- and it can be possessed whole,
as the uniform answer of an entire neighborhood, every direction at once, with one
guarantee covering them all. The directional probe is Gateaux's derivative; the
whole-neighborhood object is Frechet's; and they are not the same thing --- a
function can answer every direction separately and still refuse to have a uniform
neighborhood answer. Their ledger row is earned here, and on exactly the
distinction the record uses: the coupling is second-variation content of the
whole-neighborhood kind, and what a box shows is values --- results of proddings
already performed --- never the manner in which values would vary. To a reader of
the finished box, a directional answer and a neighborhood answer look identical:
the box records what happened, and the difference between the two derivatives is a
fact about what would have happened. The two are indistinguishable to a box-reader.
That is not a rhetorical flourish; it is the content of the theorem.

And so the fence, proved: a box-reading cannot separate the first variation from
the second. Whatever is tabulated from finished states, however cleverly, admits
two readings --- one in which the tabulated behavior is rate, one in which it is
curvature --- and nothing in the box breaks the tie, because breaking it requires
paying elaboration, which is precisely what a box-reading is defined not to do.
The coupling lives on the far side of that tie. A method that cannot tell the
second variation from the first cannot extract what only the second contains. The
proof, the record notes, is choice-free --- no oracle anywhere in it, in keeping
with the absence Chapter 2 made content --- and the reader should hear what the
theorem does not say. It does not say the coupling cannot be derived; universal
impossibilities are cheap to claim and dear to prove, and this book deals in
neither. It says this method cannot: the box-reading, precisely defined, provably
cannot. The narrowness is the honesty. A fence that claimed one inch beyond its
proof would be the very contraband this book exists to refuse.

\begin{fence}
No reading of the finished box derives the coupling. Inspection of the machine's
completed states --- taken without paying elaboration --- provably cannot separate
the first variation from the second, and the coupling lives in the second. The
claim is exactly this narrow: this method cannot. The theorem is choice-free.
\end{fence}

Now the reason this fence is the book's proudest, and it is not modesty --- it is
vindication. Every chapter of labor in this argument --- the heartbeats counted,
the proximity slip measured at its posts, the mediant walk run to the fencepost,
the cold rebuilds, the immunity theorems --- all of it would be mere procedure,
one aesthetic among possible aesthetics, if the coupling could have been read off
the box. The theorem says it could not. The number is not in the box; it is in
the walking --- in the paid elaboration the box-reading is forbidden by
definition to spend. Measurement was forced, not chosen. This book did not prefer
the laborious path; it proved the effortless one closed. The fence does not
diminish the argument. It explains why the argument had to exist.

One theorem, then, and the chapter's remaining fences are of a different kind:
lines the argument could cross and will not, marked as chosen honesty rather
than proved impossibility. Those lines --- what this instrument speaks of, and
what it borrows --- are next.

\section{The lines}
% Node-map: F2 at FULL strength (the misappropriation fence, DISPLAYED: the electron
% and nothing else; the export list stated; WHY the line sits there -- an exported
% number arrives with no earn-site behind it, its receipt does not travel); the
% chosen-vs-proved taxonomy carried from S1's closer (a chosen line honestly kept
% binds as hard as a theorem -- it binds differently); F3 in ONE referencing
% paragraph (Ch5's fence stands as built; situated, never re-entered). No new
% citation rows (Gauss waits for S3 by design). Gate: turn 531 (Ch7 S2, ~900 words).

The last section's fence was a theorem: the crossing it forbids is impossible, and
no strength of desire re-opens it. The fences of this section are of the other
kind named in that section's closing taxonomy --- lines the argument could cross,
in the trivial sense that sentences claiming more could always be typed --- and
will not. The reader should not hear "chosen" as "weaker." A chosen line honestly
kept binds as hard as a theorem; it binds differently. A theorem cannot be
violated; a vow can only be kept, and the keeping --- visible, checkable, page by
page --- is the content. This book has kept one vow through seven chapters, and
this section displays it.

The argument concerns the electron, and nothing else. Here is what that sentence
forbids, stated plainly enough that no reader can misappropriate this book by
accident. No other particle inherits anything from these pages: not a number, not
a bracket, not a method-claim. The nucleus inherits nothing; no force beyond the
one coupling measured here is spoken for; no scale above or below the electron
model's own is implicated; no cosmology, no unification, no theory of anything
else is advanced, seeded, or winked at. The coupling bracket is the electron
model's bracket. The balance reading claims kinship of architecture with a shed
in Clapham and nothing more --- Chapter 5 fenced it so. The minus twenty-one
belongs to the tower of self-description that produced it. If a reader leaves
these pages carrying a number into any other physics, the number was stolen, not
given.

And the line sits exactly where it does for the ground's own reason, not for
modesty's. Every number in this argument names the text it was read off --- that
has been the law since Chapter 2 --- and the texts are, all of them, the electron
model's own. An exported number would arrive at its new home with no earn-site
behind it: the first question any receiving discipline should ask --- off which
text was this read? --- has no answer in the new domain, because the receipt does
not travel. The number does not change in transit; its earning does not survive
the crossing. Any honest bridge from these readings to any other physics would
therefore be new work in its entirety: new texts, new measurements, new
earn-sites, new fences, built and named on the far shore. Such a bridge may be
possible. This book builds none, attempts none, and implies none.

\begin{fence}
This argument speaks of the electron and of nothing else. No other particle, no
nucleus, no force, no scale, no cosmology inherits a number, a bracket, or a
method-claim from these pages. Any bridge from these readings to any other
physics would be new work --- built, named, and fenced on its own ground --- and
this book builds none. A number carried out of this book is a number without a
receipt.
\end{fence}

One more word about whom this fence is written for, because it is not chiefly the
skeptic. The skeptic attacks the argument where it stands, and the falsifier
table welcomes them. The fence guards against the argument's friends: the
enthusiast who would praise this book by extending it --- to the proton by
analogy, to gravity by enthusiasm, to everything by induction --- does it more
damage than any attacker, because the extension claims precisely what the
argument never earned, and the debt lands on this book's name. So the standing
invitation extends here too, in the same words the ground used: if this book
anywhere borrows the electron's prestige for a claim about anything else, that
is a defect, and the reader is asked to report it as one.

Among the chosen lines stands one already built, and it is situated here rather
than rebuilt. Chapter 5's frame fence --- the unmeasured temperature, the
borrowed picture of the accelerating observer, the taper's discipline at hard
cuts --- stands exactly as displayed there, and this section adds nothing inside
it. It belongs to this family: a line drawn where proof is unavailable and
honesty is; the shape stated, the crossing refused, the borrowing credited on
instrument-not-authority terms. One theorem, two vows: that is the fence
inventory so far.

What remains for the fences chapter is not another line but a practice: what
this argument does when it meets a gap it can neither close nor ignore. There
are three honest fates, the book has used all three, and one of them once
demoted a great mathematician's rule from load-bearing to side-diagnostic
without erasing him from the ledger. The fates, the open bracket, and that
demotion --- the fence chapter's close --- are next.

\section{The fates and the bracket}
% Node-map: F4 as PRACTICE, shown not defined (the three fates with the book's own
% record as worked examples: CLOSE = the immunity theorems answering the incident;
% RESTATE/MOVE = the circular count refused and re-earned, and GAUSS'S DEMOTION as
% the fullest worked example -- his row lands; the honest ledger keeps its demotions);
% FENCE = the temperature, referenced not re-entered. F5 the OPEN BRACKET as the
% family it is (principle stated once, instances referenced; forcing a point = how a
% fake number is made). Chapter closes on the fence inventory (one theorem, two vows,
% three fates, one discipline) handing Ch9 the door.
% Gate: turn 532 (Ch7 S3, ~1000 words, chapter closer).

Chapter 1 named them in passing, in the section on how to read: a gap in an argument
has three honest fates --- close it, restate it so it no longer arises, or fence it.
The reader has now watched all three work across six chapters without being told
which was which. The fence chapter's last duty is to show the practice plainly,
because what an argument does with its gaps is a better signature than what it does
with its successes.

The first fate is closure, and its worked example is the book's oldest wound. The
incident of Chapter 3 opened a gap with teeth: if one tick could move a count, what
protected the readings? The gap was closed --- by proof, the three for-every-load
cancellation theorems --- and the reader should notice what closure leaves behind:
nothing. No mark, no caveat, no standing worry; a closed gap leaves a theorem where
the anxiety was, and the argument walks on it thereafter as on any floor.

The second fate is restatement --- the move --- and it has two worked examples of
different sizes. The small one the reader knows: the circular count of Chapter 2
could not be closed (no proof could make the cycle's period a lawful count as it
stood), so it was restated --- the three re-earned from the side-by-side text ---
after which the original worry no longer arises: the question "may the cycle's
period be the count?" dissolved, chapters later, into "the period is a theorem of
the earned count." The gap was not defeated; it was made meaningless, which is the
move fate's signature.

The larger example deserves its history told whole, because it involves the demotion
of a great mathematician's rule, and how an argument demotes is as telling as how it
proves. Early in this argument's history, a rule of Gauss's --- the three-point
rule, reading a quantity through three stations at once --- was borrowed and stood
load-bearing near the bracket's foundations. It worked; Gauss's rules do. But when
the mediant walk was built --- the machine cutting with its own knife, every cut
earned --- the borrowed rule was superseded: the walk did natively, with receipts,
what the rule did by import. The three-point rule was not deleted. It was demoted,
in the open: moved out of the load path, kept in the apparatus, and marked for what
it now is --- a side-diagnostic, consulted but never leaned on. And the ledger keeps
the row. An honest ledger includes its demotions, because the history of what was
borrowed --- including what was borrowed, used well, and set down --- is part of the
attribution; Gauss's row stands in this book's back matter not despite the demotion
but with it, labeled. The alternative, quiet deletion, would falsify the argument's
history, and an argument that falsifies its history has no standing to ask trust for
its measurements. A demotion done in the open is the move fate working exactly as
designed, and it is the last citation this book's chapters owe.

The third fate is the fence, and its worked example the reader has already stood
before: the temperature. That gap can be neither closed --- the device contains no
temperature, and no proof can conjure one --- nor restated, because the landing's
shape genuinely calls for what the instrument genuinely lacks. So it was fenced,
and Chapter 5's fence stands as built: the shape stated, the denominator refused,
the fence itself the content. Three fates, three exhibits, all from the book's own
record: proof where proof could reach, restatement where the ground allowed
re-earning, and a fence where honesty was the only instrument left.

One discipline remains, and it crowns the family because every reading in the book
lives inside it. Measurement, under this book's rules, ends in an interval whose
walls are proved ordered --- the reader has seen the lids of the flagship, the
corridor of the polygons, the standing fences of the walk, each pair with its
ordering machine-checked cold. The principle behind the instances, stated once: the
strain of not-knowing relaxes into the interval's width, and the width is the
honesty. Forcing a point out of an interval is how a fake number is manufactured
--- precision asserted beyond the walls actually built, the counterfeit's exact
signature --- and this book's intervals therefore have two possible fates and only
two: they die honest deaths, a wall failing its cold check and the reading dying
with it, or they stand as intervals. There is no third outcome, and the value near
137.011 remains on its final mention what it was on its first: a display between
walls, asserted never.

The fence inventory is complete, and it is short enough to carry: one theorem ---
the box cannot yield the coupling, and the labor was therefore necessary; two vows
--- the electron only, and the frame's borrowed picture held at arm's length; three
fates --- close, move, fence, each with its exhibit on the record; and one
discipline --- the walls stay ordered and the brackets stay open. An argument that
knows this precisely where it stops has earned something in return: the right to
say what it expects next, and to be held to it. The last chapter opens that door.
